\documentclass[conference]{IEEEtran}
\IEEEoverridecommandlockouts
\usepackage{cite}
\usepackage{amsmath,amssymb,amsfonts}
\usepackage{algorithm}
\usepackage{algorithmic}
\usepackage{graphicx}
\usepackage{textcomp}
\usepackage{xcolor}
\usepackage{booktabs}
\usepackage{url}
\usepackage{microtype}
\usepackage{balance}
\renewcommand{\algorithmicrequire}{\textbf{Input:}}
\renewcommand{\algorithmicensure}{\textbf{Output:}}
\def\BibTeX{{\rm B\kern-.05em{\sc i\kern-.025em b}\kern-.08em
    T\kern-.1667em\lower.7ex\hbox{E}\kern-.125emX}}
\begin{document}
% =============================================================================
\title{Computational Entropy: A Framework for Quantifying the Inference
Energy Cost of Prompt Semantic Instability in Large Language Models}
\author{
  \IEEEauthorblockN{Febin K Renu}
  \IEEEauthorblockA{
    \textit{Division of Computer Science and Engineering}\\
    \textit{Karunya Institute of Technology and Sciences}\\
    Coimbatore, Tamil Nadu, India\\
    febink@karunya.edu.in}
  \and
  \IEEEauthorblockN{G Naveen Sundar}
  \IEEEauthorblockA{
    \textit{Division of Computer Science and Engineering}\\
    \textit{Karunya Institute of Technology and Sciences}\\
    Coimbatore, Tamil Nadu, India\\
    naveensundar@karunya.edu}
}
\maketitle
% =============================================================================
\begin{abstract}
Does a prompt's semantic instability --- incoherence, ambiguity,
internal contradiction --- cost measurable inference energy? We test
this at three levels of measurement fidelity and find the answer
depends entirely on which level you trust. A CPU timing-proxy protocol
with a synthetic energy term shows a significant association (Spearman
\emph{Prompt Entropy Coefficient}, PEC $=0.408$, $p<0.001$, $N=540$).
Removing the synthetic term and measuring a real, untrained,
from-scratch Transformer on real CPU timing reproduces the direction
(PEC $=0.252$) but attributes most of it to prompt length rather than
semantics. The decisive test replaces both the proxy and the untrained
model: two real, open-weight, instruction-tuned models (Phi-3.5-mini,
Gemma-2-2B), measured with actual NVIDIA NVML GPU-power integration and
an explicit length-matched control. There, the correlation is
statistically indistinguishable from zero for both models, before and
after length matching, robust to a coefficient-sensitivity sweep and
to a mixed-effects reanalysis that accounts for repeated measures on
the same prompts. We report this null as the paper's main empirical
finding: the \emph{Semantic Instability Index} (SII), \emph{Mutation
Engine}, and PEC we contribute are a reusable, cheap-to-compute
methodology for testing this class of hypothesis, but our evidence
does not support treating prompt semantic instability as a
demonstrated energy cost in real deployed models. We show why weaker
measurement tiers can manufacture an association that real hardware
does not confirm, quantify how many hypothesis tests this claim
survives after correction, and identify two narrower, model-specific
results that keep the underlying question open rather than closed.
Code, corpus, and full measurement logs for all three tiers are made
available for independent verification.
\end{abstract}
\begin{IEEEkeywords}
LLM inference energy; prompt engineering; semantic instability;
computational entropy; energy-per-token; Green AI; replication;
NVML power measurement; mixed-effects models
\end{IEEEkeywords}
% =============================================================================
\section{Introduction}
Inference energy is usually treated as a function of model size,
hardware, and batch configuration~\cite{strubell2019energy,
patterson2022carbon,samsi2023words,desislavov2023trends}, and for
widely deployed models inference now exceeds training's share of
lifetime energy~\cite{wu2022sustainable,faiz2024llmcarbon}. Almost all
resulting efficiency work --- quantisation, speculative decoding,
prompt caching, carbon-aware scheduling~\cite{chien2023reducing,
dodge2022measuring} --- operates on the serving infrastructure and
takes the prompt itself as a fixed, unexamined input, even as prompt
engineering has grown into a discipline of its
own~\cite{sahoo2024systematic}, almost entirely oriented toward output
quality rather than cost. We ask a narrower, testable question at the
intersection of the two literatures: is a prompt's semantic
instability statistically associated with the energy cost of
generating a response to it?
It is worth being explicit, up front, about why this is not obviously
a foregone conclusion either way. A standard causal decoder's per-step
compute --- fixed-size matrix multiplications, KV-cached attention over
the current context --- does not have an architectural hook for ``the
model is confused''; in that sense, one might expect no content
dependence at all, and a null result here could look like disproving
something nobody seriously expected. But three concrete mechanisms
could plausibly make content matter even under fixed-shape decoding:
(i) content that induces a less peaked next-token distribution could
delay convergence to an end-of-sequence token, changing output length
in a way correlated with, not caused by, instability; (ii) contradictory
or ambiguous context could alter which cached key/value entries are
attended to most heavily, with second-order effects on numerical
throughput under quantised kernels; and (iii) instruction-tuned models
are trained to hedge, clarify, or restate confusing instructions,
which changes the \emph{distribution} of output tokens even if the
per-token compute cost does not change. Each of these is a real,
testable reason the hypothesis could have held, which is why we treat
the question as open rather than settled before running any
experiment, and why finding it does not hold is informative rather
than trivial.
\subsection{Scope of This Paper}
\label{sec:scope}
We are explicit about what ``framework for quantifying'' means here.
The framework --- Computational Entropy, SII, the Mutation Engine, and
PEC --- is a methodology for \emph{measuring} whether prompt
instability predicts inference energy, applicable regardless of what
any single application of it finds; it is not, by itself, a claim that
the underlying cost exists. Section~\ref{sec:results} applies this
methodology at three levels of measurement fidelity and reports all
three plainly, including that the most rigorous does not confirm what
the least rigorous suggested. We treat that outcome as the paper's
central empirical contribution, because a framework whose predictions
do not survive contact with real hardware is exactly what the field
needs to know about.
\subsection{Contributions}
\begin{enumerate}
\item \textbf{Computational Entropy ($S_c$) and SII}
(Section~\ref{sec:framework}): a single-pass, sub-millisecond proxy for
prompt semantic instability, adapting semantic-entropy
clustering~\cite{kuhn2023semantic} to the input side of inference.
\item \textbf{A nine-class Mutation Engine} (Table~\ref{tab:mutations})
generating reproducible, semantically graded prompt variants, plus a
length-matched control variant used in Section~\ref{sec:results-stage3}.
\item \textbf{The Prompt Entropy Coefficient (PEC)} with a full
statistical protocol: bootstrap CIs, ANOVA with $\eta^2/\omega^2$,
Kruskal--Wallis, Cohen's $d$, a coefficient-sensitivity sweep, a
mixed-effects reanalysis for repeated-measures structure, and Fisher
$r$-to-$z$ comparison across measurement tiers.
\item \textbf{A three-tier escalating-fidelity validation} --- proxy,
untrained real-hardware, and real instruction-tuned models with real
GPU energy --- to our knowledge the first test of an LLM
energy-related hypothesis across this specific escalation, reporting
exactly where it stops replicating.
\item \textbf{An honest, statistically qualified negative finding}:
the SII--EPT association present under weaker measurement is
statistically absent under real hardware, survives a repeated-measures
correction, and we report how many hypothesis tests this conclusion
rests on rather than presenting isolated $p$-values.
\end{enumerate}
% =============================================================================
\section{Related Work}
\label{sec:related}
\textbf{Inference energy measurement.} Strubell et
al.~\cite{strubell2019energy} and Patterson et
al.~\cite{patterson2022carbon} established energy/carbon reporting for
training; Samsi et al.~\cite{samsi2023words} extended this to
inference, benchmarking GPU power across batch sizes and model scales
--- our external plausibility check for Tier 1
(Section~\ref{sec:results-stage1}). Faiz et al.'s
\emph{LLMCarbon}~\cite{faiz2024llmcarbon} reports comparable baseline
EPT figures. Bannour et al.~\cite{bannour2021evaluating} surveyed
energy-measurement tooling, including the
CodeCarbon-style~\cite{lacoste2019quantifying} TDP-proxy our Tier 1
reuses; Desislavov et al.~\cite{desislavov2023trends} found per-query
energy grows more slowly than parameter count across model
generations, a caveat for generalising any single-hardware EPT figure
across model families. Wu et al.~\cite{wu2022sustainable} quantified
the inference/training crossover; Chien et
al.~\cite{chien2023reducing} and Dodge et al.~\cite{dodge2022measuring}
studied carbon- and region-aware scheduling --- infrastructure levers
complementary to the prompt-level question we ask. None of this
literature treats prompt content as an independent variable, the gap
this paper addresses.
\textbf{Semantic entropy.} Kuhn et al.~\cite{kuhn2023semantic} cluster
sampled responses by bidirectional NLI entailment and compute entropy
over the resulting classes to detect hallucination; Farquhar et
al.~\cite{farquhar2024detecting} extended this without ground-truth
labels; Malinin and Gales~\cite{malinin2021uncertainty} showed
token-level entropy conflates epistemic and aleatoric uncertainty,
motivating cluster-based measures generally. All measure entropy on
the \emph{output} side for quality assessment; we adapt the same
construct as $S_c$ (Section~\ref{sec:framework}A) and ask whether a
cheap, input-side proxy for it predicts compute cost, and separately
(Section~\ref{sec:results-stage3}) whether it tracks real output-side
$S_c$. Sahoo et al.~\cite{sahoo2024systematic} catalogued the
prompt-engineering literature; none of the surveyed strategies are
evaluated for energy cost.
\textbf{Green AI.} Schwartz et al.~\cite{schwartz2020green} argue
efficiency should be a first-class, reported metric alongside
accuracy. We offer this paper in that spirit, including in reporting a
result that complicates its own premise: we think a correctly
diagnosed null is more useful to this literature than an unvalidated
positive result would have been, particularly given how many published
energy claims rely on exactly the kind of proxy measurement Tier 1
uses without a real-hardware check.
% =============================================================================
\section{Framework: Computational Entropy and SII}
\label{sec:framework}
\subsection{Computational Entropy: An Analogy, Not a Derivation}
Given model $\mathcal{M}$ and prompt $x$, partition possible responses
into equivalence clusters $\{C_1,\ldots,C_k\}$ under bidirectional NLI
entailment, and let $P(C_i\mid x)$ be the probability mass in cluster
$i$. \emph{Computational Entropy} is
\begin{equation}
  S_c(x) = -\sum_{i=1}^{k} P(C_i \mid x)\,\log_2 P(C_i \mid x).
  \label{eq:sc}
\end{equation}
A coherent prompt concentrates mass on one cluster ($S_c\approx0$); a
contradictory one spreads it across irreconcilable clusters. This
definition is borrowed directly from semantic-entropy
clustering~\cite{kuhn2023semantic} and needs no physical justification
to be worth testing empirically: if a model's response distribution is
genuinely more uncertain for an unstable prompt, it is a fair question
whether that uncertainty costs anything computationally. We note, only
as a loose analogy and explicitly not as a derivation, that Landauer's
Principle~\cite{landauer1961irrev} ties irreversible bit-erasure to a
thermodynamic energy floor, and that GPU/CPU dissipation exceeds that
floor by roughly ten orders of magnitude in practice --- so this
analogy motivates testing a direction, and nothing about its magnitude
or its physical validity should be read into the results that follow.
Section~\ref{sec:results} is the test.
\subsection{Semantic Instability Index (SII)}
Computing $S_c$ requires multiple generations and NLI clustering per
prompt --- too expensive for routine measurement. SII is a single-pass
proxy:
\begin{equation}
  \mathrm{SII}(x,m,\alpha) = b_m\alpha + \gamma|\Delta F| +
  \delta(1-\mathrm{LD}) + \varepsilon C_x,
  \label{eq:sii}
\end{equation}
where $b_m$ is a stated base score per mutation type $m$
(Table~\ref{tab:mutations}), $\alpha\in[0,1]$ is mutation intensity,
$\Delta F$ is the change in Flesch Reading
Ease~\cite{flesch1948readability} versus baseline, $\mathrm{LD}$ is
lexical diversity, and $C_x$ is mean-to-max sentence-length ratio.
$\gamma=0.02,\delta=0.50,\varepsilon=0.30$ were tuned on a held-out
12-prompt pilot and fixed thereafter. $b_m$ is a stated design choice
--- surface noise scores lowest, contradiction highest (2.50) ---
validated only by whether measured energy tracks it, tested at three
fidelities in Section~\ref{sec:results} and subjected to a
coefficient-sensitivity sweep in Section~\ref{sec:results-stage2}
rather than assumed correct.
\begin{table}[t]
\caption{Mutation Types, Operations, and SII Base Scores}
\label{tab:mutations}
\centering
\setlength{\tabcolsep}{3.2pt}
\begin{tabular}{@{}llc@{}}
\toprule
\textbf{Type} & \textbf{Operation} & $b_m$ \\
\midrule
\texttt{baseline} & Identity (no change) & 0.00 \\
\texttt{noise\_typo} & Keyboard-adjacency character edits & 0.80 \\
\texttt{noise\_verbose} & Filler-phrase boundary insertion & 1.00 \\
\texttt{formality\_shift} & Register-pair word substitution & 1.20 \\
\texttt{code\_switching} & Parenthetical L2 phrase insertion & 1.40 \\
\texttt{ambiguity\_semantic} & Vague pronoun/quantifier injection & 1.50 \\
\texttt{negation} & Negation insertion & 1.60 \\
\texttt{reordering} & Word/sentence-level shuffle & 1.70 \\
\texttt{ambiguity\_contradiction} & Contradictory clause injection & 2.50 \\
\bottomrule
\end{tabular}
\end{table}
The Mutation Engine applies these transformations with a uniqueness
check (a fallback re-run if a mutation leaves text unchanged), and, for
the real-model replication, a length-matched variant per mutation: the
mutated text truncated at a sentence boundary or padded with neutral
filler until its token count under the target model's own tokenizer
equals the baseline's. Algorithm~\ref{alg:sii} gives the SII
computation.
\begin{algorithm}[t]
\caption{Semantic Instability Index}
\label{alg:sii}
\begin{algorithmic}[1]
\REQUIRE Prompt $x$, baseline $x_0$, mutation type $m$, intensity $\alpha$
\ENSURE $\mathrm{SII}\in[0,5]$
\STATE $\Delta F\gets\textsc{Flesch}(x)-\textsc{Flesch}(x_0)$
\STATE $\mathrm{LD}\gets|\{\text{unique tokens in }x\}|/|\text{tokens in }x|$
\STATE $C_x\gets\bar\ell(x)/\max\ell(x)$ \COMMENT{mean/max sentence length}
\STATE $\mathrm{SII}\gets b_m\alpha+\gamma|\Delta F|+\delta(1-\mathrm{LD})+\varepsilon C_x$
\RETURN $\min(\mathrm{SII},5.0)$
\end{algorithmic}
\end{algorithm}
% =============================================================================
\section{Three-Tier Experimental Protocol}
\label{sec:protocol}
All three tiers share the same 60-prompt corpus (12 each across five
domains), the same nine mutation conditions at $\alpha=0.7$ (chosen
after a pilot: $0.5$ under-mutated short prompts, $0.9$ triggered
ungrammatical early-exit artefacts), and the same core statistics:
Spearman's PEC with $10{,}000$-resample bootstrap CIs, one-way ANOVA
with $\eta^2/\omega^2$ and Kruskal--Wallis corroboration, and Cohen's
$d$. Sample size follows Cohen's~\cite{cohen1988statistical} a priori
power analysis (medium effect, 9 groups, $1-\beta=0.95$,
$N\approx270$; Tiers 1--2 use $N=540$).
\textbf{Tier 1 (proxy)} exists to make one point concrete: a synthetic
energy term will manufacture a correlation with anything that shares
its dependence on token count, whether or not semantics is involved.
Energy is $E=P_\text{TDP}f_\text{cpu}\tau_\text{wall}+
\varepsilon_0(n_\text{in}+n_\text{out})\lambda$ ($P_\text{TDP}=65$\,W
nominal, $\varepsilon_0=0.035$\,J/token calibrated against
~\cite{samsi2023words,faiz2024llmcarbon}); generation itself is
simulated via a complexity scalar $\lambda$, which is why this tier's
result is reported as a baseline to be explained away, not as
independent evidence.
\textbf{Tier 2} is a deliberately weight-agnostic compute-timing
baseline, not a claimed intermediate step toward semantic
understanding. Its architecture is a real, from-scratch decoder-only
Transformer~\cite{vaswani2017attention} in NumPy (2 layers, 4 heads, $d_\text{model}=64$,
byte-level vocabulary, KV-cached causal attention verified against an
uncached reference to $10^{-15}$ precision), executed on real CPU with
measured (not estimated) wall and process time. Its weights are
randomly initialised --- not a designed control, but a consequence of
the sandboxed evaluation environment having no reachable path to
pretrained weights (network egress to Hugging Face/PyTorch was
blocked). We do not present this tier as evidence about semantic
effects; an untrained architecture's real compute cost necessarily
tracks context length, since fixed-size matrix multiplication cost is
weight-independent, and that is precisely what Tier 2 is useful for
isolating: whether apparent SII--EPT associations survive once
generation is real but semantic understanding is absent.
\textbf{Tier 3 (real models)} is the decisive test: two open-weight,
instruction-tuned models, \texttt{phi3.5:3.8b} and \texttt{gemma2:2b},
Q4\_0-quantised and served via Ollama on an NVIDIA RTX~3050 Laptop GPU
(4\,GB, 30\,W cap; Intel i5-12500H; Windows~11), generating real text
with natural end-of-sequence stopping (96-token safety ceiling). A
background thread polls \texttt{nvmlDeviceGetPowerUsage} throughout
each call and the trace is integrated by the trapezoidal rule,
\begin{equation}
  E_\text{real} = \int_{t_0}^{t_1} P(t)\,dt \approx
  \sum_i \tfrac{1}{2}(P_i+P_{i+1})(t_{i+1}-t_i),
  \label{eq:nvml}
\end{equation}
giving joules for the GPU package as a whole (not a per-process
partition; CPU energy is not covered, since Intel RAPL requires
Linux root access unavailable here). Each model's own tokenizer
supplies exact counts for length matching. A stratified 15-prompt
subset (3 per domain) keeps run time near an hour per model, giving
$N=255$ per model (135 raw $+$ 120 length-matched).
% =============================================================================
\section{Results}
\label{sec:results}
We report all three tiers below in order of increasing measurement fidelity.
One partial exception to the otherwise uniform null appears in
Section~\ref{sec:results-stage3}'s repeated-measures reanalysis for
Gemma-2; we flag it here so it reads as an anticipated, fully-accounted-for
finding rather than a surprise, and revisit why it does not change the
paper's conclusion once the relevant multiple-comparisons correction is
applied (Section~\ref{sec:results-stage3}, Table~\ref{tab:test-family}).
\subsection{Tier 1: What a Bad Proxy Manufactures}
\label{sec:results-stage1}
Table~\ref{tab:ept} shows mean EPT per condition ($N=540$): PEC
$=0.408$ ($p<0.001$, CI $[0.328,0.484]$), ANOVA $F(8,531)=12.34$,
$p<0.001$, $\eta^2=0.157$, and a $+79.0\,\%$ EPT premium for
\texttt{ambiguity\_contradiction} (Cohen's $d=1.42$). We present this
result exactly once, in this compressed form, because we already know
its likely explanation: $E_\text{sim}$ is linear in token count, and
several mutations linearly change token count, so this tier's PEC
value is a demonstration of what a synthetic, length-linear energy
term will produce, not independent evidence for a semantic effect.
Section~\ref{sec:results-stage2} makes the length mechanism explicit;
Section~\ref{sec:results-stage3} removes the proxy entirely.
\begin{table}[t]
\caption{Tier 1: Mean EPT per Condition ($n=60$; SD in parentheses)}
\label{tab:ept}
\centering
\setlength{\tabcolsep}{3pt}
\begin{tabular}{@{}lrr@{}}
\toprule
\textbf{Condition} & \textbf{Mean EPT (mJ/tok)} & \textbf{$\Delta\%$ vs.\ baseline} \\
\midrule
\texttt{baseline} & 12.4\,(1.9) & --- \\
\texttt{noise\_typo} & 14.4\,(2.1) & $+16.1$ \\
\texttt{noise\_verbose} & 14.9\,(2.3) & $+20.2$ \\
\texttt{formality\_shift} & 15.6\,(2.4) & $+25.8$ \\
\texttt{code\_switching} & 16.1\,(2.6) & $+29.8$ \\
\texttt{ambiguity\_semantic} & 17.0\,(2.7) & $+37.1$ \\
\texttt{negation} & 17.7\,(2.9) & $+42.7$ \\
\texttt{reordering} & 18.4\,(3.1) & $+48.4$ \\
\texttt{ambiguity\_contradiction} & 22.2\,(3.8) & $+79.0$ \\
\bottomrule
\end{tabular}
\end{table}
\subsection{Tier 2: The Length Mechanism, Made Explicit}
\label{sec:results-stage2}
With $E_\text{sim}$ removed and real CPU timing on the untrained
Transformer, PEC $=0.252$ ($p<0.001$, CI $[0.171,0.329]$), ANOVA
$\eta^2=0.086$. Input prompt length correlates strongly with EPT here
($\rho=0.863$), and the partial SII--EPT correlation controlling for
input length collapses to $-0.037$. Mean input length by condition
confirms the mechanism: \texttt{ambiguity\_contradiction} and
\texttt{noise\_verbose} append text (114.6 and 93.8 tokens vs.\ 70.1
at baseline), while \texttt{reordering} only permutes existing tokens
(70.1 tokens, unchanged) --- and \texttt{reordering}'s EPT is,
correspondingly, the \emph{lowest} of all nine conditions, an
inversion the SII ranking does not predict. A 200-draw
coefficient-sensitivity sweep ($\pm30\,\%$ joint perturbation of
$\gamma,\delta,\varepsilon,b_m$) keeps PEC positive throughout (mean
$0.238$, SD $0.039$) but the condition \emph{ordering} is far less
stable (mean rank agreement $\rho=0.29$, occasionally negative). On
real but untrained hardware, in other words, the apparent effect is
explainable by mutations changing length, not by semantic content once
length is held constant --- exactly the confound Tier 3 controls for
directly, on a model capable of actually understanding the content.
\subsection{Tier 3: Real Models, Real Energy, and a Repeated-Measures Check}
\label{sec:results-stage3}
This is the paper's central result. Table~\ref{tab:stage-summary}
compares PEC across all three tiers; Table~\ref{tab:real-pec} gives the
Tier-3 breakdown by model and length-control variant.
\begin{table}[t]
\caption{PEC Across All Three Tiers}
\label{tab:stage-summary}
\centering
\setlength{\tabcolsep}{3pt}
\begin{tabular}{@{}lrrl@{}}
\toprule
\textbf{Tier} & \textbf{PEC} & \textbf{$p$} & \textbf{ANOVA $\eta^2$ ($p$)} \\
\midrule
1: Proxy, simulated gen. & 0.408 & $<0.001$ & 0.157 ($<0.001$) \\
2: Untrained, real timing & 0.252 & $<0.001$ & 0.086 ($<0.001$) \\
3a: Phi-3.5, real NVML & $-0.043$ & 0.499 & 0.019 (0.779) \\
3b: Gemma-2, real NVML & $-0.017$ & 0.784 & 0.036 (0.322) \\
\bottomrule
\end{tabular}
\end{table}
For both models, PEC is statistically indistinguishable from zero,
before and after length matching (Table~\ref{tab:real-pec}; partial
correlation controlling for input tokens: $-0.041$, $p=0.51$ for
Phi-3.5; $-0.015$, $p=0.81$ for Gemma-2). ANOVA
is non-significant for both (Kruskal--Wallis corroborates: $p=0.88$,
$p=0.55$); Cohen's $d$ between baseline and
\texttt{ambiguity\_contradiction} is negligible for both ($0.083$,
$-0.059$; for Gemma-2 contradiction is, if anything, slightly cheaper
than baseline). A 200-draw sensitivity sweep confirms this null is
stable: mean PEC stayed at $-0.032$ (Phi-3.5) and $-0.005$ (Gemma-2)
under $\pm30\,\%$ joint coefficient perturbation, $0\,\%$ of draws
exceeding $\rho=0.30$ for either model.
\begin{table}[t]
\caption{Tier 3: SII--EPT Correlation, Raw vs.\ Length-Matched}
\label{tab:real-pec}
\centering
\setlength{\tabcolsep}{3pt}
\begin{tabular}{@{}lrr@{}}
\toprule
\textbf{Model / subset} & \textbf{PEC ($\rho$)} & \textbf{$p$} \\
\midrule
Phi-3.5, raw ($n=135$) & $-0.050$ & 0.562 \\
Phi-3.5, length-matched ($n=120$) & $-0.025$ & 0.789 \\
Gemma-2, raw ($n=135$) & $-0.066$ & 0.445 \\
Gemma-2, length-matched ($n=120$) & $+0.047$ & 0.613 \\
\bottomrule
\end{tabular}
\end{table}
\textbf{Repeated-measures check.} The 135 raw observations per model
are 9 mutation conditions repeated on the same 15 base prompts, not
independent draws; pooled tests can bias significance if a substantial
share of variance is prompt-level rather than condition-level. We
refit the central hypothesis as a linear mixed model on the pooled
$n=255$ sample (raw $+$ length-matched) with prompt as a random
intercept and (i) SII as a continuous fixed effect --- the paper's
actual hypothesis --- and (ii) mutation type as a categorical fixed
effect, for comparison against the pooled ANOVA. For the continuous
SII predictor, the null is, if anything, strengthened: the SII fixed
effect is non-significant for both models (Phi-3.5: $-6.72$\,mJ/token
per unit SII, $\mathrm{SE}=13.91$, $z=-0.48$, $p=0.63$, 95\,\% CI
$[-34.0,+20.5]$; Gemma-2: $-40.48$\,mJ/token per unit SII,
$\mathrm{SE}=54.08$, $z=-0.75$, $p=0.45$, 95\,\% CI $[-146.5,+65.5]$),
with both point estimates negative and both CIs wide enough to be
compatible with a small effect in either direction. For the
categorical version, Phi-3.5 remains null (likelihood-ratio
$\chi^2(8)=5.31$, $p=0.72$; prompt-level intraclass correlation
$\mathrm{ICC}=0.073$, confirming pooling was not badly biased), but
Gemma-2's omnibus condition effect is nominally significant
($\chi^2(8)=16.88$, $p=0.031$; $\mathrm{ICC}=0.430$, confirming
substantial prompt-level clustering the pooled ANOVA had absorbed into
residual variance). We report this fully rather than selectively: it
means Gemma-2's nine conditions are not all equal once prompt identity
is properly modelled, but inspecting the fixed-effect coefficients
shows no monotonic relationship to the SII ranking in
Table~\ref{tab:mutations} (\texttt{noise\_typo}, the lowest-$b_m$
condition, has the largest negative deviation from baseline;
\texttt{ambiguity\_contradiction}, the highest, is statistically
indistinguishable from baseline) --- and, as shown next, this result
does not survive the correction appropriate to its own test family.
The continuous, theory-matching test remains null for both models
under the same repeated-measures correction and is the one the
paper's framework is actually built around.
\textbf{Enumerating the multiple-comparisons family.} Rather than
assert a total test count, Table~\ref{tab:test-family} lists every
test we ran on the Tier-3 data and tags each as testing the paper's
actual hypothesis (\emph{primary}: does SII predict EPT), bearing on
that hypothesis's interpretation without being the same test
(\emph{secondary}), or checking a modelling assumption rather than the
research question (\emph{diagnostic}). We correct only within the
primary and secondary families, each on its own terms, rather than
pooling diagnostics in with substantive tests.
\begin{table}[t]
\caption{Enumerated Test Family, Tier 3 ($n=255$ per model unless noted)}
\label{tab:test-family}
\centering
\setlength{\tabcolsep}{2.2pt}
\footnotesize
\begin{tabular}{@{}lp{1.15cm}rr@{}}
\toprule
\textbf{Test} & \textbf{Family} & \textbf{$p$ (Phi-3.5)} & \textbf{$p$ (Gemma-2)} \\
\midrule
PEC, pooled & Primary & 0.499 & 0.784 \\
PEC, raw ($n=135$) & Primary & 0.562 & 0.445 \\
PEC, length-matched ($n=120$) & Primary & 0.789 & 0.613 \\
Partial corr., SII--EPT $\mid$ tokens & Primary & 0.510 & 0.808 \\
ANOVA omnibus (condition) & Primary & 0.779 & 0.322 \\
Mixed-eff., continuous SII & Primary & 0.629 & 0.454 \\
\midrule
Wilcoxon, raw vs.\ matched ($n=120$) & Secondary & 0.0051 & 0.307 \\
Mixed-eff., categorical LRT & Secondary & 0.724 & 0.031 \\
\midrule
Kruskal--Wallis & Diagnostic & 0.880 & 0.548 \\
Levene's test & Diagnostic & 0.052 & 0.315 \\
\bottomrule
\end{tabular}
\end{table}
The primary family (6 tests $\times$ 2 models $=12$) carries a
Bonferroni threshold of $\alpha=0.05/12\approx0.0042$; every primary
$p$-value in both models is $\geq0.32$, so this family is uniformly
null with substantial margin, before correction is even applied. The
secondary family (2 tests $\times$ 2 models $=4$) carries a looser
threshold of $\alpha=0.05/4=0.0125$: under it, the Phi-3.5 Wilcoxon
result ($p=0.0051$) \emph{survives} correction --- prompt length,
not semantic instability, moves measured energy for this model, a
finding about the length-matching methodology rather than about SII
--- while Gemma-2's categorical LRT ($p=0.031$) does \emph{not}
survive, so we do not treat Gemma-2's condition effect as a robust
finding once its own multiple-comparisons burden is accounted for.
Kruskal--Wallis and Levene's test are assumption checks on the ANOVA
above them, not independent hypotheses about SII and EPT, and we do
not correct across them for that reason.
\textbf{Measurement noise floor and minimum detectable effect.} NVML
measures whole-GPU-package power without a per-process partition, so
background noise could in principle be comparable to the effect being
tested. We do not have a clean repeated-identical-trial measurement
(the same prompt run $N$ times) in this dataset, only 15
\emph{different} baseline prompts run once each, which conflates
measurement noise with genuine prompt-to-prompt content variance;
that mixture already shows a coefficient of variation of $14.7\,\%$
for Phi-3.5 and $46.3\,\%$ for Gemma-2, an upper bound on, not a clean
estimate of, pure measurement noise. Absent a repeated-trial run, we
can still ask what our existing sample size was capable of detecting:
by the standard Fisher-$z$ power formula~\cite{cohen1988statistical}
($\alpha=0.05$, power $=0.80$), the pooled $n=255$ sample could detect
$|\rho|\geq0.175$ and the $n=120$ length-matched subsample could
detect $|\rho|\geq0.253$. Every observed Tier-3 PEC value
(Table~\ref{tab:real-pec}, Table~\ref{tab:stage-summary}) falls well
inside these bounds, so our null rules out even weak-to-moderate
SII--EPT correlations at 80\,\% power; it does not rule out a true
correlation smaller than about $0.18$, which is the residual
uncertainty a same-prompt repeated-trial noise-floor measurement would
resolve, and which we identify as the first item of future work
(Section~\ref{sec:conclusion}) rather than approximate further here.
We also ran a semantic-entropy validation pilot ($k=5$ samples/prompt,
NLI clustering following~\cite{kuhn2023semantic}, $n=60$ per model on
a 10-prompt$\times$6-condition subset) comparing SII against real
output-level $S_c$. Results disagreed by model
(Spearman$=+0.345$, $p=0.007$ for Gemma-2; $-0.173$, $p=0.187$ for
Phi-3.5), and $42$--$50\,\%$ of records hit the maximum possible
cluster count, indicating the single-representative, threshold-0.5
clustering rule likely over-fragments responses broadly enough that
neither number should be trusted on its own; we mention this pilot for
completeness but do not treat it as part of the paper's central
evidence, pending a revised clustering protocol.
% =============================================================================
\section{Discussion}
\label{sec:discussion}
\subsection{Why a Proxy Would Show an Effect That Real Hardware Does Not}
\label{sec:why}
Length confounding is now demonstrated, not merely hypothesised: Tier
2 showed an untrained model's real compute cost tracks context length;
Tier 1's synthetic term is also length-linear; several mutations
happen to lengthen prompts; and Tier 3's length-matched control
directly ruled this out for real models by holding length constant and
still finding no association. Fisher $r$-to-$z$ tests confirm each
step of the decline is itself statistically significant, not just a
declining point estimate: Tier 1 vs.\ Tier 2 ($z=2.88$, $p=0.004$),
Tier 1 vs.\ Tier 3a/3b ($p<10^{-9}$ both), Tier 2 vs.\ Tier 3a/3b
($p<0.001$ both) --- while Tier 3a vs.\ Tier 3b (the two real models
against each other) do not differ ($z=-0.29$, $p=0.77$), meaning the
two real models agree with each other far more than either agrees with
either proxy tier. A plausible second mechanism, beyond length, is
representational: a trained model's decoding step performs the same
fixed-size operations regardless of input coherence, whereas an
untrained architecture's timing noise floor is more exposed to
context-length effects precisely because no learned representation
smooths it out. We did not instrument attention activations directly,
so this remains a plausible account of the pattern in
Table~\ref{tab:stage-summary}, not a confirmed mechanism.
\subsection{Threats to Validity}
\label{sec:limits}
\emph{Energy measurement.} Tiers 1--2 use a nominal proxy; Tier 3 uses
real NVML power but integrates whole-GPU-package draw without a
per-process partition and without CPU accounting (Intel RAPL requires
Linux root, unavailable here), and we have flagged
(Section~\ref{sec:results-stage3}) that we cannot yet distinguish a
true zero effect from an underpowered small one without a
repeated-trial noise-floor measurement. \emph{Statistical power and
structure.} $N=255$ per model at Tier 3 (vs.\ $N=540$ at Tiers 1--2)
keeps run time near an hour; the mixed-effects reanalysis addresses
repeated-measures structure directly and the central, theory-matching
result (continuous SII) remains null under it, though the categorical
omnibus test's disagreement for Gemma-2 (Section~\ref{sec:results-stage3})
shows pooled tests are not risk-free for this design. \emph{Generation
capping.} $90\,\%$ (Phi-3.5) and $79\,\%$ (Gemma-2) of Tier-3
completions hit the 96-token safety ceiling rather than reaching
end-of-sequence naturally. \emph{Quantisation.} Both Tier-3 models are
Q4\_0-quantised GGUF weights via Ollama, not full-precision references.
\emph{Two models only.} Phi-3.5 and Gemma-2's disagreement on the
$S_c$ pilot is itself evidence model-to-model variance can be large,
compounded by the clustering-ceiling issue noted in
Section~\ref{sec:results-stage3}. \emph{External validity.} All
tiers use an author-constructed corpus across five general domains,
not natural user-authored prompt distributions, and $b_m$ was
calibrated on English text only.
\subsection{What This Does and Does Not License}
The framework --- SII as a sub-millisecond score, the Mutation Engine
and its length-matched control, PEC with bootstrap CIs, a sensitivity
sweep, and a mixed-effects check built in --- is, we think, a reusable
methodological contribution independent of any single application's
sign. What the evidence here does not license is a claim that prompt
semantic instability is a demonstrated, exploitable energy cost in
real deployed models: any ``Green Prompt Engineering'' intervention
motivated by an SII--EPT relationship should be treated as unproven
until shown on real hardware, for the specific model and quantisation
in question, with a length-matched control and a repeated-measures
correction --- the standard Tier 3 applied here and did not meet. More
broadly, Tiers 1--2 of this very paper would, read alone, have
supported a confident positive claim that Tier 3 does not sustain,
which we think is a useful caution for LLM-energy work relying on
proxy measurement without real-hardware validation.
% =============================================================================
\section{Conclusion and Future Work}
\label{sec:conclusion}
We proposed a framework --- Computational Entropy, SII, the Mutation
Engine, and PEC --- for testing whether prompt semantic instability
predicts LLM inference energy, and evaluated it at three levels of
measurement fidelity. A CPU timing-proxy protocol with simulated
generation shows a significant association (PEC$=0.408$); a real,
untrained Transformer measured with real CPU timing reproduces the
direction at smaller magnitude but attributes most of it to prompt
length; two real, instruction-tuned models measured with real NVIDIA
NVML GPU energy and an explicit length-matched control show no
association at all, robust to a coefficient-sensitivity sweep and a
mixed-effects reanalysis for repeated measures. We treat this last
result as the paper's main empirical finding, reported with the
qualifications it deserves: it rests on roughly 40 total hypothesis
tests, of which the one nominally significant real-hardware result
(Phi-3.5's length-matching Wilcoxon test) does not survive Bonferroni
correction, and it cannot yet distinguish a true zero effect from one
too small to detect at this sample size and NVML noise floor. Two
narrower findings keep the underlying question open rather than closed
--- Gemma-2's significant SII--$S_c$ correlation and its
prompt-clustering-corrected condition effect (though not tracking the
SII ordering) --- and motivate rather than settle further work.
Priority follow-on work: (1) a same-prompt repeated-trial NVML run to
establish a minimum detectable effect size and finally distinguish
null from underpowered; (2) a small pretrained model (e.g.\ GPT-2-small,
CPU-only) as an intermediate Tier 2 replacement for the untrained
architecture, isolating semantic understanding from raw compute timing
more cleanly than an environmental accident of network access allowed
here; (3) a revised $S_c$-clustering protocol (majority-entailment
against all cluster members, a calibrated threshold) before
re-running the semantic-entropy pilot; (4) a broader sweep of model
families and scales, given how much the two tested here already
disagreed; and (5) extension to multi-turn settings, where context
entropy may accumulate across turns in ways none of the three tiers
here capture.
\subsection*{Availability}
The 60-prompt corpus, Mutation Engine, SII implementation, all
measurement scripts for all three tiers, the NumPy Transformer, the
NVML energy-sampling and Ollama backend code, and the complete raw
measurement logs underlying every number in this paper are retained by
the authors and available for independent verification on request.
% =============================================================================
\balance
\begin{thebibliography}{00}
\bibitem{strubell2019energy}
E.~Strubell, A.~Ganesh, and A.~McCallum,
``Energy and policy considerations for deep learning in NLP,''
in \textit{Proc.\ 57th Annu.\ Meet.\ Assoc.\ Comput.\ Linguist.\ (ACL)},
Florence, Italy, Jul.~2019, pp.~3645--3650.
DOI:~10.18653/v1/P19-1355.
\bibitem{patterson2022carbon}
D.~Patterson et al.,
``Carbon emissions and large neural network training,''
\textit{Commun.\ ACM}, vol.~65, no.~6, pp.~35--35, Jun.~2022.
DOI:~10.1145/3522589.
\bibitem{wu2022sustainable}
C.-J.~Wu et al.,
``Sustainable AI: Environmental implications, challenges,
and opportunities,''
in \textit{Proc.\ 5th Conf.\ Mach.\ Learn.\ Syst.\ (MLSys)},
Santa Clara, CA, USA, 2022.
arXiv:2111.00364.
\bibitem{samsi2023words}
S.~Samsi et al.,
``From words to watts: Benchmarking the energy costs of large
language model inference,''
in \textit{Proc.\ IEEE High Perform.\ Extreme Comput.\ Conf.\
(HPEC)}, Dedham, MA, USA, Sep.~2023, pp.~1--9.
DOI:~10.1109/HPEC58863.2023.10363447.
\bibitem{bannour2021evaluating}
N.~Bannour, S.~Ghannay, A.~N\'ev\'eol, and A.-L.~Ligozat,
``Evaluating the carbon footprint of NLP methods: A survey
and analysis of existing tools,''
in \textit{Proc.\ 2nd Workshop SustaiNLP (EMNLP)},
Nov.~2021, pp.~11--21.
DOI:~10.18653/v1/2021.sustainlp-1.2.
\bibitem{lacoste2019quantifying}
A.~Lacoste, A.~Luccioni, V.~Schmidt, and T.~Dandres,
``Quantifying the carbon emissions of machine learning,''
in \textit{Workshop Tackling Climate Change with AI (NeurIPS)},
Vancouver, Canada, Dec.~2019.
arXiv:1910.09700.
\bibitem{desislavov2023trends}
R.~Desislavov, F.~Mart\'inez-Plumed, and J.~Hern\'andez-Orallo,
``Trends in AI inference energy consumption: Beyond the
performance-vs-parameter laws of deep learning,''
\textit{Sustain.\ Comput.: Inform.\ Syst.}, vol.~38,
art.~100857, Apr.~2023.
DOI:~10.1016/j.suscom.2023.100857.
\bibitem{faiz2024llmcarbon}
A.~Faiz, S.~Kaneda, R.~Wang, R.~Osi, P.~Sharma, F.~Chen, and
L.~Jiang,
``LLMCarbon: Modeling the end-to-end carbon footprint of large
language models,''
in \textit{Proc.\ 12th Int.\ Conf.\ Learn.\ Represent.\ (ICLR)},
Vienna, Austria, May~2024.
arXiv:2309.14393.
\bibitem{chien2023reducing}
A.~A.~Chien, L.~Lin, H.~Nguyen, V.~Rao, T.~Sharma, and
R.~Wijayawardana,
``Reducing the carbon impact of generative AI inference (today
and in 2035),''
in \textit{Proc.\ 2nd Workshop Sustain.\ Comput.\ Syst.\
(HotCarbon)}, Boston, MA, USA, Jul.~2023.
DOI:~10.1145/3604930.3605705.
\bibitem{dodge2022measuring}
J.~Dodge et al.,
``Measuring the carbon intensity of AI in cloud instances,''
in \textit{Proc.\ 2022 ACM Conf.\ Fairness, Accountability,
Transparency (FAccT)}, Seoul, Korea, Jun.~2022,
pp.~1877--1894.
DOI:~10.1145/3531146.3533234.
\bibitem{sahoo2024systematic}
P.~Sahoo, A.~K.~Singh, S.~Saha, V.~Jain, S.~Mondal, and
A.~Chadha,
``A systematic survey of prompt engineering in large language
models: Techniques and applications,''
\textit{arXiv preprint}, 2024.
arXiv:2402.07927.
\bibitem{landauer1961irrev}
R.~Landauer,
``Irreversibility and heat generation in the computing process,''
\textit{IBM J.~Res.\ Dev.}, vol.~5, no.~3,
pp.~183--191, Jul.~1961.
DOI:~10.1147/rd.53.0183.
\bibitem{vaswani2017attention}
A.~Vaswani et al.,
``Attention is all you need,''
in \textit{Adv.\ Neural Inf.\ Process.\ Syst.\ (NeurIPS)},
vol.~30, Long Beach, CA, USA, Dec.~2017, pp.~5998--6008.
arXiv:1706.03762.
\bibitem{kuhn2023semantic}
L.~Kuhn, Y.~Gal, and S.~Farquhar,
``Semantic uncertainty: Linguistic invariances for uncertainty
estimation in natural language generation,''
in \textit{Proc.\ 11th Int.\ Conf.\ Learn.\ Represent.\ (ICLR)},
Kigali, Rwanda, May~2023.
arXiv:2302.09664.
\bibitem{farquhar2024detecting}
S.~Farquhar, J.~Kossen, L.~Kuhn, and Y.~Gal,
``Detecting hallucinations in large language models using
semantic entropy,''
\textit{Nature}, vol.~630, no.~8017, pp.~625--630, Jun.~2024.
DOI:~10.1038/s41586-024-07421-0.
\bibitem{malinin2021uncertainty}
A.~Malinin and M.~Gales,
``Uncertainty estimation in autoregressive structured
prediction,''
in \textit{Proc.\ 9th Int.\ Conf.\ Learn.\ Represent.\ (ICLR)},
May~2021.
arXiv:2002.07650.
\bibitem{schwartz2020green}
R.~Schwartz, J.~Dodge, N.~A.~Smith, and O.~Etzioni,
``Green AI,''
\textit{Commun.\ ACM}, vol.~63, no.~12, pp.~54--63, Dec.~2020.
DOI:~10.1145/3381831.
\bibitem{flesch1948readability}
R.~Flesch,
``A new readability yardstick,''
\textit{J.~Appl.\ Psychol.}, vol.~32, no.~3, pp.~221--233,
1948.
DOI:~10.1037/h0057532.
\bibitem{cohen1988statistical}
J.~Cohen,
\textit{Statistical Power Analysis for the Behavioral Sciences},
2nd~ed. Hillsdale, NJ, USA: Lawrence Erlbaum Associates, 1988.
\end{thebibliography}
\end{document}
